\documentclass[conference]{IEEEtran}

\usepackage{graphicx}
\usepackage{amsmath,amssymb}
\usepackage{float}
\usepackage{booktabs}
\usepackage{multirow}
\usepackage{xcolor}

\begin{document}

\title{
    \textbf{Benchmark Report: 505.mcf\_r} \\
    \text{Register Spill Analysis on RISC-V via gem5}
}

\author{
    \IEEEauthorblockN{Lütfullah Burak Kaya}
    \IEEEauthorblockA{
        Ozyegin University \\
        burak.kaya.50711@ozu.edu.tr
    }
    \and
    \IEEEauthorblockN{Ismail Akturk}
    \IEEEauthorblockA{
        Ozyegin University \\
        ismail.akturk@ozyegin.edu.tr
    }
}

\newcommand{\LongDate}{%
    \number\day\space
    \ifcase\month
    \or January\or February\or March\or April\or May\or June%
    \or July\or August\or September\or October\or November\or December%
    \fi
    \space \number\year
}

\twocolumn[
\begin{@twocolumnfalse}
\maketitle
\begin{center}{\small \LongDate}\end{center}
\vspace{1em}
\end{@twocolumnfalse}
]

% ================================================

\begin{abstract}
This report presents the register spill characterization of the
\texttt{505.mcf\_r} benchmark from SPEC CPU2017, executed on a RISC-V
(RV64GC) target using the gem5 TimingSimpleCPU simulator.
A custom SpillDetector was integrated into gem5 to count store-then-load
patterns on the stack within a user-defined Region of Interest (ROI).
The test dataset (5,985 nodes, 84,449 arcs) was used.
Two simulation runs were attempted: the first failed with an
out-of-memory error; the second, with 4\,GB of simulated memory,
completed successfully. Results show a spill rate of 1.83\%
within the ROI, corresponding to 446 million detected spill events
across 24.4 billion simulated instructions.
\end{abstract}

% =========================================================
\section{Benchmark Overview}

\texttt{505.mcf\_r} is the SPEC CPU2017 rate variant of MCF, a
minimum-cost flow optimizer used for public transit scheduling.
The algorithm constructs a transportation network as a graph and
solves it using the successive shortest path / simplex method.
The benchmark is memory-intensive and pointer-heavy, performing
large numbers of graph traversals.

The RISC-V binary (\texttt{mcf\_r.riscv}) was compiled with gem5 ROI
markers inserted around the main optimization call in \texttt{mcf.c}:

\begin{itemize}
    \item \texttt{m5\_work\_begin(0,0)} — inserted before \texttt{global\_opt()}
    \item \texttt{m5\_work\_end(0,0)}   — inserted after  \texttt{global\_opt()}
\end{itemize}

Only spills occurring inside \texttt{global\_opt()} are counted
toward the reported statistics.

% =========================================================
\section{Input Datasets}

MCF reads a single graph file (\texttt{inp.in}) whose first line
encodes the problem size as \texttt{<nodes> <arcs>}.
Table~\ref{tab:inputs} summarises all four standard datasets.

\begin{table}[H]
\centering
\caption{MCF Input Dataset Sizes}
\label{tab:inputs}
\begin{tabular}{lrrl}
\toprule
\textbf{Dataset} & \textbf{Nodes} & \textbf{Arcs} & \textbf{Lines in inp.in} \\
\midrule
test     &  5,985 &  84,449 &  90,435 \\
train    & —      & —       & 177,967 \\
refrate  & 15,000 & 117,202 & 132,203 \\
refspeed & —      & —       & 210,494 \\
\bottomrule
\end{tabular}
\end{table}

\textbf{Only the test dataset was run} in this study.
The test \texttt{inp.in} file is 1.2\,MB (90,435 lines).
Even with the smallest dataset, the simplex algorithm executes
71,002 iterations internally, making it one of the heavier
benchmarks despite its modest graph size.

% =========================================================
\section{Simulation Setup}

\subsection{gem5 Configuration}

\begin{itemize}
    \item \textbf{CPU model:}   TimingSimpleCPU
    \item \textbf{ISA:}         RISC-V RV64GC
    \item \textbf{Caches:}      enabled (default L1 I/D + L2)
    \item \textbf{Memory:}      4\,GB simulated physical RAM
    \item \textbf{Spill mode:}  summary only (\texttt{spill\_verbose=False})
\end{itemize}

\subsection{Run History}

Two simulation runs were performed on the same test input, as
summarised in Table~\ref{tab:runs}.

\begin{table}[H]
\centering
\caption{MCF Simulation Run Summary}
\label{tab:runs}
\begin{tabular}{llll}
\toprule
\textbf{Run} & \textbf{Mem} & \textbf{Outcome} & \textbf{Output} \\
\midrule
Run 1 (old binary) & 512\,MB & \textcolor{red}{OOM crash} & 37.1\,GB verbose log \\
Run 2 (new binary) & 4\,GB   & \textcolor{green!50!black}{Completed} & $\approx$300\,B summary \\
\bottomrule
\end{tabular}
\end{table}

\textbf{Run 1} used the previous verbose SpillDetector (every spill
written to disk) with the default gem5 memory of 512\,MB.
MCF's simplex algorithm dynamically expands its arc storage via
\texttt{mremapping}, eventually exhausting the simulated physical memory
at simplex iteration 57,610. The run was terminated by gem5 with
\texttt{Out of memory}; no valid ROI summary was produced.

\textbf{Run 2} used the refactored counter-only SpillDetector with
4\,GB of simulated memory. The simulation ran to completion (simplex
iteration 71,002, checksum verified) and produced a compact summary file.
All reported results below are from Run 2.

The full gem5 invocation for Run 2 was:

\begin{verbatim}
./build/RISCV/gem5.opt
  --outdir=benchmarks/mcf/m5out_test_summary
  configs/deprecated/example/se.py
  --cmd=benchmarks/mcf/mcf_r.riscv
  --options="benchmarks/mcf/inp.in"
  --cpu-type=TimingSimpleCPU
  --caches
  --mem-size=4GB
\end{verbatim}

% =========================================================
\section{Simulation Results — Test Input}

\subsection{Pre-ROI Context}

Before entering \texttt{global\_opt()}, the binary reads and constructs
the network graph. Table~\ref{tab:preroi} shows global counters at the
moment \texttt{m5\_work\_begin} fired.

\begin{table}[H]
\centering
\caption{Global Counters at ROI Entry}
\label{tab:preroi}
\begin{tabular}{lr}
\toprule
\textbf{Metric} & \textbf{Value} \\
\midrule
Total instructions (pre-ROI) & 194,085,337 \\
Total spills detected (pre-ROI) & 19,671,795 \\
\bottomrule
\end{tabular}
\end{table}

\subsection{ROI Spill Statistics}

Table~\ref{tab:roi} presents the spill statistics collected exclusively
within the ROI (\texttt{global\_opt()} call, test input).

\begin{table}[H]
\centering
\caption{ROI Spill Statistics — 505.mcf\_r (test input)}
\label{tab:roi}
\begin{tabular}{lr}
\toprule
\textbf{Metric} & \textbf{Value} \\
\midrule
ROI instructions         & 24,426,982,048 \\
ROI stores               &  1,300,618,520 \\
ROI loads                &  8,113,038,194 \\
\midrule
\textbf{ROI spills}      & \textbf{446,113,136} \\
\textbf{Spill rate}      & \textbf{1.8263\%} \\
\bottomrule
\end{tabular}
\end{table}

\subsection{Timing}

\begin{table}[H]
\centering
\caption{Simulation Timing — Run 2 (test input, completed)}
\label{tab:timing}
\begin{tabular}{lr}
\toprule
\textbf{Metric} & \textbf{Value} \\
\midrule
Simulated time  & 147.42\,s \\
Host wall time  & 17,512\,s ($\approx$4.9\,hours) \\
\bottomrule
\end{tabular}
\end{table}

% =========================================================
\section{Analysis}

\subsection{Spill Rate Interpretation}

The measured spill rate of \textbf{1.83\%} means that approximately
1 in every 55 instructions inside the ROI is a register spill.
This is notably higher than the pre-ROI phase:

\[
\text{Pre-ROI spill rate} = \frac{19{,}671{,}795}{194{,}085{,}337} \approx 10.1\%
\]

The graph-construction phase (pre-ROI) is therefore more
register-pressure-intensive than the optimization loop itself.

\subsection{Store vs.\ Load Balance}

Inside the ROI, loads (8.11B) outnumber stores (1.30B) by a factor
of approximately 6.2$\times$. This reflects MCF's access pattern:
the simplex iterations repeatedly read graph-node fields
(\texttt{cost}, \texttt{potential}, \texttt{flow}) without rewriting
them every access, so the same spilled values are reloaded many more
times than they are written.

\subsection{Spill-to-Load Ratio}

\[
\frac{446{,}113{,}136}{8{,}113{,}038{,}194} \approx 5.50\%
\]

Roughly 1 in 18 load instructions inside the ROI is a spill reload,
indicating moderate but sustained register pressure throughout the
optimization phase.

\subsection{Memory Requirement}

Run 1 (512\,MB) crashed at simplex iteration 57,610 out of 71,002.
Run 2 (4\,GB) completed cleanly. MCF's dynamic arc expansion
(\texttt{mremapping}) requires substantially more simulated physical
memory than the default gem5 SE configuration provides.
Future runs of larger datasets (refrate/refspeed) will also
require at least \texttt{--mem-size=4GB}.

% =========================================================
\section{Conclusion}

The \texttt{505.mcf\_r} benchmark exhibits a spill rate of 1.83\%
within its core \texttt{global\_opt()} ROI on RISC-V RV64GC.
With 446 million spill events across 24.4 billion instructions,
MCF presents a higher spill burden than NAMD (1.03\%) from the same
study, consistent with its pointer-chasing, graph-traversal nature
which stresses register allocation more than floating-point pipelines.
The graph-construction phase (pre-ROI) reaches 10.1\%, making it
the highest-pressure region of the entire run.
Larger dataset runs (refrate: 15,000 nodes / 117,202 arcs) are
expected to scale the spill count roughly proportionally to the
increase in simplex iterations.

\end{document}
